\chapter{Data dalam Microservices: Saga dan Outbox}

\section*{Tujuan Pembelajaran}
\begin{itemize}
  \item Menjelaskan alasan setiap layanan memiliki datanya sendiri dan risiko basis data bersama.
  \item Merancang saga beserta langkah kompensasinya untuk satu proses lintas layanan.
  \item Menjelaskan cara pola \textit{transactional outbox} mencegah data dan \textit{event} tidak sinkron.
\end{itemize}

\section{Pendahuluan}

% TODO
% Cakupan: data pasien Medikara Health tersebar di banyak sistem dan sering tidak cocok satu sama lain; memecah sistem menjadi layanan tanpa memikirkan data justru memperparah masalah ini.

\section{Database per Service dan Basis Data Bersama}

% TODO
% Cakupan: alasan tiap layanan memiliki datanya sendiri; basis data bersama sebagai anti-pattern dan dampaknya pada perubahan dan kepemilikan; cara layanan lain memperoleh data yang bukan miliknya.
% Kasus Medikara: tabel pasien SIMRS yang ditulis oleh pendaftaran, keuangan, dan laboratorium sekaligus.

\section{Transaksi Terdistribusi dan Eventual Consistency}

% TODO
% Cakupan: mengapa satu transaksi basis data tidak bisa lagi mencakup seluruh proses; keterbatasan two-phase commit; eventual consistency dan cara menjelaskannya kepada pengguna bisnis; konsekuensi pada pengalaman pasien.

\section{Saga dan Transaksi Kompensasi}

% TODO
% Cakupan: gagasan saga dari Garcia-Molina dan Salem (1987); saga berbasis choreography dan orchestration; transaksi kompensasi dan langkah yang tidak dapat dibatalkan; pola saga dalam Richardson (2018).
% Kasus Medikara: saga pendaftaran pasien rujukan dari klinik mitra sampai klaim terkirim.

\section{Pola Transactional Outbox}

% TODO
% Cakupan: masalah dual write (data tersimpan tetapi event tidak terkirim, atau sebaliknya); cara kerja transactional outbox; kaitan dengan idempotensi consumer (Bab 9).
% Pendalaman teknis: menerapkan outbox sederhana dengan PostgreSQL dan proses pembaca yang menerbitkan event.

\section{CQRS dan Read Model untuk Pelaporan}

% TODO
% Cakupan: CQRS pada tingkat konsep; read model yang dibangun dari event untuk kebutuhan laporan; manfaatnya bagi layanan operasional; batasan kesegaran data.
% Kasus Medikara: laporan manajemen yang masih disusun dengan spreadsheet sebagai calon read model.

\section{Kepemilikan Data Pasien: Satu Pasien, Satu Rekam Medis}

% TODO
% Cakupan: identitas pasien tunggal lintas fasilitas; layanan pemilik data induk pasien dan pembaca lainnya; persetujuan pasien dan hak akses sebagai batasan desain (pelindungan data pribadi dibahas rinci di IT40813); arsitektur data tingkat enterprise disinggung singkat (dibahas rinci di IT40513).
% Kasus Medikara: strategi "Satu Pasien, Satu Rekam Medis" dan pertukaran data dengan SATUSEHAT.

\section{Aktivitas: Peta Kepemilikan Data dan Rancangan Saga}

% TODO
% Tujuan aktivitas: menetapkan pemilik setiap entitas data utama dan merancang satu saga lintas layanan beserta kompensasinya.
% Peran dalam tim: anggota non-teknis memastikan aturan bisnis dan langkah kompensasi masuk akal bagi pasien dan bagian keuangan; anggota teknis memilih pola koordinasi dan titik pemakaian outbox; seluruh anggota menguji rancangan dengan skenario gagal.
% Langkah: (1) daftar entitas data utama; (2) tetapkan layanan pemilik dan layanan pembaca; (3) pilih satu proses lintas layanan; (4) uraikan langkah saga dan kompensasinya; (5) uji dengan dua skenario gagal.
% Luaran dan tonggak capstone: C3, peta kepemilikan data dan rancangan saga.
% Berkas sources yang direncanakan: sources/chapter10/templat-peta-kepemilikan-data.xlsx, sources/chapter10/templat-rancangan-saga.md, sources/chapter10/contoh-outbox-postgresql/ (opsional)

\section{Ringkasan}

% TODO

\section{Latihan}

\begin{enumerate}
  \item Mengapa tim keuangan dan tim layanan pasien sebaiknya tidak menulis langsung ke tabel basis data yang sama?
  \item Pada saga pendaftaran rujukan, verifikasi kepesertaan BPJS Kesehatan gagal setelah jadwal dokter dikunci. Uraikan langkah kompensasinya.
  \item Laporan manajemen Medikara Health masih disusun dengan \textit{spreadsheet}. Jelaskan cara \textit{read model} menggantikan kebiasaan itu tanpa membebani layanan operasional.
\end{enumerate}
